\section{Served cards}\label{sec:hook}

Most of the time an agent does not know that the store holds something relevant, so it does not ask. The \emph{hook} addresses this. It is a program that both agent runtimes run on every prompt an agent receives, before the agent reads it, and whose output is added to the agent's context. When the prompt is mathematical and the store holds a strong match, the hook prints at most a few cards, each in one or two lines, with its identifier and a line saying how far to trust it. Otherwise it prints nothing.

\subsection{Design constraints}
The hook runs on every prompt of every agent, so its constraints are the reverse of those of the recall tool. It must be fast (it runs in the foreground); it must be silent when it has nothing good, because an irrelevant note costs every agent attention and, worse, can be obeyed; and it must never fail loudly, so every error leaves it silent. A note from the hook is a lead, not an answer: an agent that wants to rely on it opens the card with the recall tool.

\subsection{The gate}
The hook first decides whether the prompt is mathematical. A prompt passes the lexical gate only if it contains at least two distinct \emph{mathematical anchors}: words or symbols that occur in the store's mathematical vocabulary and not in a list of general words (an anchor is, for instance, a named theorem, a defined object, a symbol such as $\Var$ or $\ell^2$, or an author's name attached to a result). A prompt that passes is scored with the base profile of Section~\ref{sec:retrieval}, and the hook abstains unless the best card clears a support threshold and leads the second by a margin. The graph and symbol channels are given less weight here than in the deliberate tool, because on real prompts they matched notation more often than meaning; and of the cards the hook shows, at most one may be a \emph{campaign deposit} (a record written by an agent during a running campaign), so that the newest material cannot crowd out the settled.

\subsection{The trust line}
Every note carries one line that states what kind of thing it is and how much it has been checked: a theorem with a proof text, a concept lead, a harvested draft, an unreviewed attempt, a refuted route. The line is there because served notes can mislead. In the $61$-item reasoning test (Section~\ref{sec:tests}) a served statement helped on items whose fault was a dropped hypothesis and hurt on items whose fault was a citation, a scope or a sharpness error: a note about the cited result made a wrong use of it look legitimate. A note says that the question resembles something in the store; it does not say that the resemblance holds, and the agent is told to weigh it against what it sees.

\subsection{What the hook achieves}
The hook is cheap and it is rarely wrong, but it answers little by itself. On the $100$ store questions it spoke on $15$ and was silent on $85$, including all $60$ questions that had been reworded to share no words with their card; its gate read those as non-mathematical, and its search had no shared words to start from. On the $15$ where it spoke, the full stack answered $11$, the recall tool alone $9$, text search $10$, and the hook's notes alone $7$. The hook is best understood as a reminder that the store exists and a pointer to where to start; the work of recall is done by the deliberate tool. We return in Section~\ref{sec:limits} to the possibility of a hook that speaks more often with the same precision.
